% V033 首页包含主教学图；保留四项结论
\begin{statementbox}
\begin{center}
\begin{tikzpicture}[x=1cm,y=1cm,font=\small,>=Stealth,
  every node/.style={inner sep=3pt},line cap=round,line join=round]
\path[use as bounding box] (-1,-1.6) rectangle (11,5.5);
\node[font=\small\bfseries,text=CStatement] at (5,5.1) {同一条对角线，把平行四边形分成两半};

\coordinate (O) at (2,0); \coordinate (A) at (6,1);
\coordinate (B) at (3,3); \coordinate (D) at (7,4);
\fill[CStatement!17] (O)--(A)--(B)--cycle;
\fill[CExample!15] (A)--(D)--(B)--cycle;
\draw[thick] (O)--(A)--(D)--(B)--cycle;
\draw[very thick,CProof] (A)--(B);
\draw[->,very thick,CStatement] (O)--(A);
\draw[->,very thick,CExample] (O)--(B);
\node[below left] at (O) {$O$}; \node[below right] at (A) {$A(x_1,y_1)$};
\node[above left] at (B) {$B(x_2,y_2)$};
\node[above right] at (D) {$D=A+B$};
\node at (3.65,1.3) {$S$}; \node at (5.35,2.7) {$S$};

\node[align=center] at (5,-1.15) {$S=\frac12|x_1y_2-x_2y_1|$，两片全等，整体面积为 $2S$。};
\end{tikzpicture}
\end{center}

\textbf{\textcolor{CStatement}{条件与结论}}
设 $O(0,0)$、$A(x_1,y_1)$、$B(x_2,y_2)$ 为任意平面点，允许负坐标、零坐标及共线退化。

\textbf{一、原点三角形面积：}
\[
\boxed{S_{\triangle OAB}=\frac12|x_1y_2-x_2y_1|.}
\]
$x_1y_2-x_2y_1$ 是有向三角形面积的两倍；普通面积必须取绝对值。

\textbf{二、平行四边形面积：} 补点 $D=(x_1+x_2,y_1+y_2)$，则
\[
\boxed{S_{OADB}=|x_1y_2-x_2y_1|=2S_{\triangle OAB}.}
\]
退化共线时两种面积均为零。

\textbf{三、任意三点面积：} 对 $P(x_1,y_1)$、$Q(x_2,y_2)$、$R(x_3,y_3)$，
将 $P$ 平移到原点，得到
\[
\boxed{S_{\triangle PQR}=\frac12\left|
(x_2-x_1)(y_3-y_1)-(x_3-x_1)(y_2-y_1)\right|.}
\]
\textbf{四、共线判定：}
\[
O,A,B\text{ 共线}\iff x_1y_2-x_2y_1=0\iff S_{\triangle OAB}=0.
\]
\textbf{学习入口：} 初中底高面积与相似三角形即可推导；过 $A$、过 $B$ 均可作水平或竖直辅助线，四种证明见方法二至五。
\end{statementbox}
